% BEGIN THERMAL_R03_SYNC_18
\section{Evaluation of the constants and quantitative comparison}
\label{sec:ii-contraste-cuantitativo}

The preceding constructions allow three operations to be separated:
producing a quantity, expressing its coordinate in a chart, and comparing that
coordinate with a reference. The APP--TRIT--TPK state preserves sheets,
orientation, quotient, carry, and memory before these readings. Its
nonadic prolongation determines the correlated coordinates
\(\pi,\varphi,e,\alpha\); the multiplicative norm over the local quotient and regional return
produce the action sections; the marked \(90/120\) incidence and
conjugate channels produce the Barbero functional; angular
transport subsequently determines the CKM coordinates. No entry from the external
tables enters the choice of these operations.

Here we collect the evaluations of the preceding constructions and their
comparison references. The digits of the internal outputs express
computational precision, not experimental uncertainty. For a value
\(x\) and a nonzero reference \(x_0\), both in the same chart, we use
\[
 \Delta x=x-x_0,\qquad r_x=\frac{x-x_0}{x_0}.
\]
When the reference has asymmetric marginal uncertainties
\(u_-,u_+\), we also define the descriptive residual
\begin{equation}
 z_x=\begin{cases}
 \Delta x/u_+,&\Delta x\geq0,\\
 \Delta x/u_-,&\Delta x<0.
 \end{cases}
 \label{eq:ii-residuo-marginal}
\end{equation}
This normalization does not assume independence between rows and does not
by itself constitute a joint statistical test.

\subsection{Reduced Planck constant and action per turn: two orientations}

In the chart \(\mathcal U_S\mapsto\mathrm{J\,s}\), the sections already
constructed are
\begin{align}
 \hbar_{\rm pre}&=H_5\,10^{-34}\,\mathrm{J\,s},&
 H_5&=1.0545718176461544696\ldots,\\
 \hbar_{\rm ret}&=\eta_{\rm ret}\,10^{-34}\,\mathrm{J\,s},&
 \eta_{\rm ret}&=1.0545718169150390611\ldots,\\
 h_a&=2\pi\hbar_a,&a&\in\{\mathrm{pre},\mathrm{ret}\}.
 \label{eq:ii-cuatro-coordenadas-planck}
\end{align}
The decade comes from \(\Sigma_H=\varphi\alpha^{16}\); the difference between
the two sections comes from regional return. Multiplication by \(2\pi\)
changes action per radian to action per turn, without introducing another
orientation or altering the relative residual.

The SI reference is
\[
 h_{\rm SI}=6.62607015\times10^{-34}\,\mathrm{J\,s},\qquad
 \hbar_{\rm SI}=\frac{h_{\rm SI}}{2\pi}
 =1.0545718176461563912\ldots\times10^{-34}\,\mathrm{J\,s}.
\]
Both expressions are exact. The infinite expansion of \(\hbar_{\rm SI}\)
is not a measurement uncertainty. Fixing the numerical value of
\(h\) belongs to the definition of the SI in effect since 2019
\cite{bipm2018}; comparison takes place after the HMT construction.

\begin{table}[htbp]
\centering\small
\setlength{\tabcolsep}{5pt}
\begin{tabular}{@{}llr@{}}
\toprule
Quantity & Coordinate in \(10^{-34}\,\mathrm{J\,s}\) & \(r_x\) relative to SI\\
\midrule
\(\hbar_{\rm pre}\) & \(1.0545718176461544696\) & \(-1.822221\times10^{-15}\)\\
\(\hbar_{\rm ret}\) & \(1.0545718169150390611\) & \(-6.932836\times10^{-10}\)\\
\(h_{\rm pre}\) & \(6.62607014999998793\) & \(-1.822221\times10^{-15}\)\\
\(h_{\rm ret}\) & \(6.62607014540625433\) & \(-6.932836\times10^{-10}\)\\
\bottomrule
\end{tabular}
\caption{Action before and after return. The final digits are
evaluations of the resulting coordinates; the residuals are not experimental
errors in \(h\) or \(\hbar\).}
\label{tab:ii-planck-comparacion}
\end{table}

\begin{proposition}[Invariance of the comparison under a complete turn]
In a common chart,
\[
 \frac{h_a-h_{\rm SI}}{h_{\rm SI}}
 =\frac{\hbar_a-\hbar_{\rm SI}}{\hbar_{\rm SI}}.
\]
\end{proposition}
\begin{proof}
The two terms of the difference in the numerator on the left contain the common factor
\(2\pi\), which cancels against that in the denominator.
\end{proof}
The pre/ret distinction is quantitatively relevant: replacing one with
the other changes the comparison by several orders of magnitude. The regional
corrector must remain visible and must not be confused with numerical rounding.

\subsection{Barbero--Immirzi: internal value and comparison between constructions}

The functional evaluation is reconstructed from its three contributions:
\begin{align*}
 \Delta S_{90}&=2.95169439297457621139847987861\times10^{-4},\\
 \Delta S_{120}&=1.96835112502951720093738706217\times10^{-5},\\
 \pi AC^*/180&=0.270485322934140078465586458790\ldots,\\
 \gamma_{\rm HMT}&=\pi AC^*/180+12\Delta S_{90}-\Delta S_{120}\\
 &=0.27400767269445927474725526077398041940\ldots.
\end{align*}
This composition preserves the oriented weight \(12:-1\), prior to its
decimal evaluation. The stated summands reconstruct the final
coordinate with an absolute residual below \(3\times10^{-31}\).

The Ghosh--Mitra counting is considered as a bibliographic comparison,
not a measurement of the parameter. Their equation (7), written with
\(\lambda=2\pi\gamma_{\rm GM}\), is
\begin{equation}
 \sum_{j\in\{1/2,1,3/2,\ldots\}}(2j+1)
 e^{-2\pi\gamma_{\rm GM}\sqrt{j(j+1)}}=1.
 \label{eq:ii-gm-comparador}
\end{equation}
The reference is A.~Ghosh and P.~Mitra,
\emph{An improved estimate of black hole entropy in the quantum geometry
approach}, \emph{Physics Letters B} \textbf{616} (2005), 114--117,
\href{https://arxiv.org/abs/gr-qc/0411035}{arXiv:gr-qc/0411035},
equations (7)--(8). The root is evaluated here for that specific equation;
puncture counting is not identified with the HMT generator.

\begin{table}[htbp]
\centering\small
\begin{tabular}{@{}lll@{}}
\toprule
Construction & Dimensionless value & Nature\\
\midrule
Functional \(\Gamma_{6\to8}\) & \(0.2740076726944593\) & HMT incidence and return\\
Root of \eqref{eq:ii-gm-comparador} & \(0.2740668582328300\) & Declared bibliographic counting\\
\midrule
\(\gamma_{\rm HMT}-\gamma_{\rm GM}\) & \(-5.918553837\times10^{-5}\) & Difference between constructions\\
\(\gamma_{\rm HMT}/\gamma_{\rm GM}-1\) & \(-2.159529202\times10^{-4}\) & Relative difference\\
\bottomrule
\end{tabular}
\caption{Barbero--Immirzi comparison in the positive orientation. No
experimental uncertainty or statistical significance is assigned to this table.}
\label{tab:ii-barbero-comparacion}
\end{table}

The external root evaluation is reproducible without arbitrarily
truncating its tail. With \(n=2j\), \(N=128\), and
\(r=e^{-0.27\pi}\), for \(\gamma\geq0.27\),
\[
 \sum_{n>N}(n+1)e^{-\pi\gamma\sqrt{n(n+2)}}
 \leq \frac{r^{N+1}[(N+2)-(N+1)r]}{(1-r)^2}<10^{-43}.
\]
Every summand decreases strictly with \(\gamma\). The endpoints
\(0.27\) and \(0.28\) enclose the root; bisection and the tail bound
control the decimals in the table. This is a check of the
bibliographic comparator, not a route for selecting \(\gamma_{\rm HMT}\).

\subsection{CKM angles, amplitudes, and CP violation}

The angular reader already constructed determines, in degrees,
\begin{equation}
 (\theta_{12},\theta_{23},\theta_{13},\delta)
 =(13.003313589509,\;2.398056157332,\;
 0.213977382244,\;65.676173123554)^\circ.
 \label{eq:ii-ckm-decimales}
\end{equation}
These values arise from \(M_{\rm CKM}(A,C^*/6,\Delta)^T\) and
\(\delta=9A\). Conversion to radians is performed once before
evaluating sines and cosines. The unitary parameterization of the angular
section gives
\[
 |V_{\rm CKM}|\simeq
 \begin{pmatrix}
 0.974350259&0.225005836&0.003734601\\
 0.224873110&0.973489279&0.041841465\\
 0.008582400&0.041121745&0.999117283
 \end{pmatrix},
\]
and the invariant quartet preserves the phase:
\[
 J=c_{12}c_{23}c_{13}^{\,2}s_{12}s_{23}s_{13}\sin\delta
 =3.1189723326632974\times10^{-5}\quad
 \text{in the angular reader evaluation.}
\]
This evaluation uses \eqref{eq:ii-jarlskog} with the coordinates of
\eqref{eq:ii-ckm-decimales}. The nine printed moduli are rounded; unitarity belongs
to the complex matrix before rounding, not to a matrix formed
solely from moduli.

The comparator is the three-generation fit in the
\href{https://pdg.lbl.gov/2025/reviews/rpp2025-rev-ckm-matrix.pdf}
{Particle Data Group, 2025 update, \emph{CKM Quark-Mixing Matrix}},
equations (12.27)--(12.28) and the adjacent value of \(J\). To avoid
altering its uncertainties through a nonlinear transformation, the table
retains the observables expressed by PDG: \(s_{ij}=\sin\theta_{ij}\),
\(\delta\) in radians, and \(J\).

\begin{table}[htbp]
\centering\small
\setlength{\tabcolsep}{4pt}
\begin{tabular}{@{}lrrr@{}}
\toprule
Observable & HMT evaluation & PDG 2025 & \(z_x\)\\
\midrule
\(s_{12}\) & \(0.225007404757\) & \(0.22501\pm0.00068\) & \(-0.003817\)\\
\(s_{23}\) & \(0.041841757010\) & \(0.04183^{+0.00079}_{-0.00069}\) & \(0.014882\)\\
\(s_{13}\) & \(0.003734601164\) & \(0.003732^{+0.000090}_{-0.000085}\) & \(0.028902\)\\
\(\delta\) (rad) & \(1.146265461116\) & \(1.147\pm0.026\) & \(-0.028251\)\\
\(10^5J\) & \(3.118972333\) & \(3.12^{+0.13}_{-0.12}\) & \(-0.008564\)\\
\bottomrule
\end{tabular}
\caption{Marginal comparison in the declared parameterization and edition.
The factor \(10^5\) jointly multiplies the value and uncertainty in the
last row; it does not change its normalized residual.}
\label{tab:ii-ckm-comparacion}
\end{table}

Residuals are computed with the coordinates before table
rounding. The PDG fit combines data and unitarity constraints;
its rows are not independent measurements. Therefore
\(\sum z_x^2\) is not summed, nor is a number of degrees of freedom assigned to this table.
The result is a reproducible marginal comparison with that edition,
not a probability measure of the model. The historical marginal deviations
from preceding calculations are not transferred to this table:
their documentation does not identify the edition and covariance of that comparison.

The PMNS sector retains the transport \(U=F_\ell^\dagger F_\nu\), the
spectral domains, and phase data developed in the angular
section. A numerical row requires specifying the frames of the record
used and the realization compared. This table does not attribute CKM numbers
to the leptonic reader or determine a new PMNS matrix.

\subsection{Boltzmann and universal gravitation: quantity, chart, and reference}

The thermal construction already recovered determines the positive transducer
\(k_B:\mathcal L_\Theta\to\mathcal L_E\) and the relation
\begin{equation}
 k_B\Theta_{{\rm clk},a}\log3
 =\frac{h_a}{108t_0}=\frac{\hbar_a}{t_P},
 \qquad t_P=\frac{54}{\pi}t_0.
 \label{eq:ii-kb-comparacion}
\end{equation}
Local information, action, and period appear before the kelvin--joule
chart. For positive bases \(u_E,u_\Theta\), if
\(k_B(u_\Theta)=b\,u_E\),
\(\Theta_{{\rm clk},a}=\theta_a u_\Theta\), and
\(E_{P,a}=\epsilon_a u_E\), the equation is equivalent to
\(b\theta_a=\epsilon_a/\log3\). A change
\(u'_E=s_Eu_E\), \(u'_\Theta=s_\Theta u_\Theta\) produces
\[
 b'=\frac{s_\Theta}{s_E}b,\qquad
 \theta'_a=\frac{\theta_a}{s_\Theta},\qquad
 \epsilon'_a=\frac{\epsilon_a}{s_E}.
\]
The marked realization of \ref{thm:ii-ter27-transducer} further
determines the coefficient through the spectral reference and the pair
information per reference--conditional energy:
\begin{equation}
 b_*=10^{-23}\left(1+380\,10^{-3}+649\,10^{-6}\right)
     =\frac{1380649}{10^{29}}=1.380649\times10^{-23}.
 \label{eq:ii-b-evaluacion-comparada}
\end{equation}
The degrees, their multiplicities and the capacity origin are retained
in \ref{sec:ii-ter27-antecedentes}; the additive characterization and
the conservative instrument are developed in
\ref{sec:ii-aditiva-instrumento}. With $H$ and the reference fixed,
$d\mathcal S/d\mathcal E=b_*\beta$ yields
$\Theta_P=(b_*\beta)^{-1}$ and $B_*(u_\Theta)=b_*u_E$.
The table therefore displays the coefficient actually evaluated,
alongside the SI reference
\[
 k_B^{\rm SI}=1.380649\times10^{-23}\,\mathrm{J\,K^{-1}},
\]
The numeral agrees exactly in this realization and in the specified
transported bases. Its thermometric proof includes the spectral measure
and both functionals: changing only their normalization changes the reader.
Identifying $u_E/u_\Theta$ with $\mathrm{J/K}$ is the dimensional
transport enabling comparison. The SI value is definitional and is
therefore assigned no experimental uncertainty.

In gravity, the circular realization fixes
\begin{equation}
 G_a=\vartheta_{108}^{\circ}\frac{c_{\rm int}^{5}t_0^2}{\hbar_a},
 \qquad
 \vartheta_{108}^{\circ}=(54/\pi)^2
 =295.4525715010569\ldots.
 \label{eq:ii-g-comparacion}
\end{equation}
The selector and reciprocity \(G_a\hbar_a=
\vartheta_{108}^{\circ}c_{\rm int}^{5}t_0^2\) are internal quantities of
the declared realization. A dimensional evaluation of \(G_a\) must
preserve, in addition to orientation \(a\), the time,
velocity, and action charts. The longitudinal and radial composition in
Section~\ref{g26:sec:rigidez-radial-en} first determines
\[
 L=\frac{5759}{23040}\alpha_{\HMT}^{16}L_*,\qquad
 \hbar_{\rm ret}=\eta_{\rm ret}S_*,\qquad
 t_0=\frac{\pi_{\HMT}L}{54c_{\rm int}}.
\]
With the same constitutive velocity section, substitution into
\eqref{eq:ii-g-comparacion} yields
\[
 G_{\rm ret}=\frac{c_{\rm int}^3L_*^2}{S_*}
 \left(\frac{5759}{23040}\right)^2
 \frac{\alpha_{\HMT}^{32}}{\eta_{\rm ret}}.
\]
Dimensional speed retains the constitutive normalization of the angular
channels. With the already generated coordinates \(A_{\rm deg}\)
and \(C^*_{\rm deg}\),
\begin{equation}
 \begin{aligned}
 q_\pm&=\exp\!\left[-\frac{\pi_{\HMT}}{180}
               (A_{\rm deg}\pm C^*_{\rm deg})\right],\\
 r_\pm&=\frac{1-q_\pm^{90}}{1-q_\pm^{120}},\qquad
 \widehat c=\frac1{r_+r_-},\qquad c_{\rm int}=c_*\widehat c.
 \end{aligned}
 \label{eq:ii-g26-velocidad-constitutiva}
\end{equation}
In compatible positive bases,
\(\varepsilon=\varepsilon_*r_+^2\), \(\mu=\mu_*r_-^2\), and
\(c_*=(\mu_*\varepsilon_*)^{-1/2}\); multiplication gives
\(c_{\rm int}=(\mu\varepsilon)^{-1/2}\). Hence
\begin{equation}
 G_{\rm ret}=\frac{c_*^3L_*^2}{S_*}
 r_N^2\alpha_{\HMT}^{32}\frac{\widehat c^{\,3}}{\eta_{\rm ret}}
 =\frac{L^2}{\hbar_{\rm ret}(\mu\varepsilon)^{3/2}}.
 \label{eq:ii-g26-gravedad-constitutiva}
\end{equation}
The evaluator uses the dimensional section \(c_{\rm int}\) declared
below and recovers \(c_*=c_{\rm int}/\widehat c\). Substitution
verifies equivalence of these expressions in the same representation:
\(\widehat c\) enters once in passing from \(c_*\) to
\(c_{\rm int}\) \cite{hmtg26counter}.

In the declared representation \(L_*=1\,\mathrm m\),
\(S_*=10^{-34}\,\mathrm{J\,s}\), and
\(c_{\rm int}=299792458\,\mathrm{m\,s^{-1}}\), the result is
\[
 L=1.616254982625126330168262501334562\ldots\times10^{-35}
 \,\mathrm m,
\]
\[
 G_{\rm ret}=6.674299659414022706367776189\ldots\times10^{-11}
 \,\mathrm{m^3\,kg^{-1}\,s^{-2}}.
\]
The comparison value \(6.67430(15)\times10^{-11}\) is retained
separately. The difference is
\(-3.40585977\ldots\times10^{-18}\,\mathrm{m^3\,kg^{-1}\,s^{-2}}\),
approximately \(-0.00227057\) standard uncertainties of that reference.
Stable computed digits describe the evaluation of the composition with
its explicit readers and bases; experimental uncertainty belongs to the
measured reference.

\begin{table}[htbp]
\centering\small
\begin{tabularx}{\linewidth}{@{}lXX@{}}
\toprule
Object & Internal evaluation & Dimensional reference\\
\midrule
\(B_*\) & \(1.380649\times10^{-23}\,u_E/u_\Theta\),
from \eqref{eq:ii-b-evaluacion-comparada} and the marked thermometric rule.
& \(1.380649\times10^{-23}\,\mathrm{J/K}\), exact in SI.\\[4pt]
\(G_{\rm ret}\) & \(6.6742996594140227\ldots\times10^{-11}\,
\mathrm{m^3/(kg\,s^2)}\), from \(c_{\rm int}^3L^2/\hbar_{\rm ret}\).
& \(6.67430(15)\times10^{-11}\,\mathrm{m^3/(kg\,s^2)}\), CODATA 2022.\\
\bottomrule
\end{tabularx}
\caption{Evaluations and references in the specified realizations.
The thermal row retains the identification of bases explained in the text.
The standard
relative uncertainty of the \(G\) reference is approximately
\(2.24743\times10^{-5}\); \(k_B\) has no definitional uncertainty.
The gravitational evaluation and its difference from the reference use
the longitudinal and radial composition declared in the text.}
\label{tab:ii-kb-g-cartas}
\end{table}

The external sources for this last table are the
\href{https://www.bipm.org/en/measurement-units/si-defining-constants}
{BIPM defining constants} and the
\href{https://physics.nist.gov/cuu/Constants/Table/allascii.txt}
{NIST CODATA 2022 table} \cite{codata2022}.
The experimental uncertainty in \(G\) is not transferred to \(h\),
\(\hbar\), or \(k_B\). Nor does an equality of digits in a chart
eliminate the need to identify the reader producing the quantity.

\subsection{Reproduction of the evaluations}

The program \texttt{verificar\_comparaciones\_II.py}, included with the
sources, receives the stored internal coordinates in files separate
from the external data. It checks the composition of the Barbero functional,
the conversion between \(h\) and \(\hbar\), the bibliographic root and its tail,
the CKM moduli, the quartet \(J\), and the residuals of those comparisons. It uses
75-digit decimal arithmetic and elementary functions evaluated after
freezing the internal outputs. Execution instructions and program versions
are included in the reproducible package.
Working precision does not increase the digits available from the
original calculation: only digits stable under its rounding are presented.
The check records the files used and their hashes, so that
a change of evaluation or external edition remains visible.

The added gravitational evaluation is reproduced by
\texttt{verificar\_reproduccion.py}, in the source directory
\texttt{reproduccion\_gravitatoria\_20260926}.
The focal set retains the tests of polar transport, metric composition,
simultaneous rigidity, phase action, and joint composition, together with
structural and angular substitution. The substitution verifiers read the
archived regional expansions and register \(K\), reconstruct the
subsequent expressions, and compare two working precisions; the
gravitational evaluations retain 71 stable significant digits.
The general identities are proved in the body; the rational checks
examine finite cases of those identities. Inputs, copy provenance,
hashes, and the scope of each program are documented in the same
directory. This execution reuses previously generated outputs and
keeps focal reproduction, complete TPK execution, and Lean
formalization distinct.

% NUMCHECK hbar_SI_coefficient=1.0545718176461563912 tol=1e-19
% NUMCHECK h_pre_coefficient=6.62607014999998793 tol=6e-18
% NUMCHECK h_ret_coefficient=6.62607014540625433 tol=6e-18
% NUMCHECK relative_pre=-1.822221e-15 tol=6e-22
% NUMCHECK relative_ret=-6.932836e-10 tol=6e-17
% NUMCHECK gamma_GM=0.2740668582328300 tol=6e-17
% NUMCHECK gamma_difference=-5.918553837e-5 tol=6e-15
% NUMCHECK gamma_relative=-2.159529202e-4 tol=6e-14
% NUMCHECK ckm.s12.value=0.225007404757 tol=6e-13
% NUMCHECK ckm.s23.value=0.041841757010 tol=6e-13
% NUMCHECK ckm.s13.value=0.003734601164 tol=6e-13
% NUMCHECK ckm.delta_rad.value=1.146265461116 tol=6e-13
% NUMCHECK ckm.s12.marginal_residual=-0.003817 tol=6e-7
% NUMCHECK ckm.s23.marginal_residual=0.014882 tol=6e-7
% NUMCHECK ckm.s13.marginal_residual=0.028902 tol=6e-7
% NUMCHECK ckm.delta_rad.marginal_residual=-0.028251 tol=6e-7
% NUMCHECK ckm.J.marginal_residual=-0.008564 tol=6e-7
% NUMCHECK gravity_selector=295.4525715010569 tol=6e-14
% NUMCHECK G_relative_reference_uncertainty=2.24743e-5 tol=6e-11
% END THERMAL_R03_SYNC_18
